% ============================================================================
% Discussion (from Arce's gBKMR_Nature.docx; to be expanded by Linda + Arce)
% ============================================================================

\section*{Discussion}

Exposure to metal mixtures has been shown to increase the risk of several chronic diseases, including diabetes and cardiovascular disease. In this study, we present a novel method, g-BKMR, to estimate the joint effect of multiple time-varying exposures over time while adjusting for time-varying confounders. Our approach allows the estimation of joint effects and interactions of time-varying environmental exposures, which is critical for understanding the complex relationships between exposures and health outcomes. Moreover, our method accounts for time-varying confounding, and could also be extended to help mitigate potential biases that may arise from unmeasured confounding factors.

The g-formula used in our method allows for complex relationships between the elements of the mixture and the confounders through the joint kernel specification in the outcome model. Beyond environmental health, our proposed method can be applied to a wide range of data with high-dimensional time-varying correlated exposures when causal interpretation is desired.

Applying g-BKMR to the Strong Heart Study, we found a harmful effect of co-exposure to six urinary metals on fasting blood glucose. Zinc was the dominant contributor to the joint mixture signal across all three visits, while selenium at Visit~1 showed a positive association that was robust to excluding participants with baseline diabetes---consistent with evidence from the Selenium and Vitamin E Cancer Prevention Trial (SELECT) that selenium supplementation elevates diabetes risk.\cite{select_trial} Cadmium at Visit~1 showed an inverse association that remained after adjustment for baseline pack-years, suggesting the pattern is not solely explained by smoking intensity. These findings highlight the importance of addressing exposure to metal mixtures in the prevention of diabetes and cardiovascular disease. Our simulation study also showed that g-BKMR outperforms other existing methods for the evaluation of environmental mixtures with repeated measures.

In conclusion, the g-BKMR method provides a powerful and flexible approach to estimate the joint effects of time-varying exposures, while accounting for time-varying confounders, and can be applied to a wide range of data. The findings from our application of g-BKMR to the Strong Heart Study underscore the importance of addressing exposure to metal mixtures in the prevention of chronic diseases. g-BKMR enables a better understanding of the complex relationship between environmental mixtures and health outcomes, paving the way for more effective public health policies and interventions.
